% !TEX root = ../main.tex
\section{Experimental Setup}
\label{sec:experiments}
\paragraph{Models.} We evaluate a ladder of eleven production reward models spanning 0.6B to 70B and the
Llama, Qwen, Gemma, and DeBERTa families, chosen to disentangle \emph{dataset-driven} from
\emph{architecture-driven} bias: it includes a same-recipe 8B/70B pair (for a controlled size comparison) and
same-data / different-architecture pairs. Results reported here are on \modelSmall{}; the larger models are
in progress. \paragraph{Metrics.} For the harm evidence (cross-marker decision design): the decision-margin
disparity per attribute with its $2^3$ decomposition (interactions and the additivity gap), decision-format
\crossinf{} (accuracy-based and threshold-free) on strong and weak records, and, as the floor, whether the
model at least penalises a decline that cites the protected attributes openly. For the mechanism layer: direct
\autoinf{}, an additivity (cosine) test for the intersection, the cross-nulling matrix of the directions of
Section~\ref{sec:method}, the placement check, a reasoning correctness$\times$conclusion $2\times2$, and linear
versus non-linear probe recoverability after \leace{}. Sample sizes are reported per arm rather than globally:
\nEvalBattery{} held-out matched pairs per battery cell and \nEvalStandpointPersuade{} (PERSUADE) /
\nEvalStandpointAsap{} (ASAP) essays for the standpoint arm; the cross-marker design's number of strong and weak
records per domain is fixed after a pilot (records used to fit a direction are never evaluated). Robustness is
assessed across encoding, template, response paraphrase and injection position.
